% [0,1]^2 在 d_inf 下的 eps 网：(a) eps=1/2，4 个中心 ((2i-1)/4,(2j-1)/4)；(b) eps=1/4，9 个中心 ((2i-1)/6,(2j-1)/6)。各示一个 d_inf 球（边长 2eps 的开正方形）。
\begin{tikzpicture}[x=3cm,y=3cm,
  c/.style={circle,fill=FigureBlue,inner sep=0pt,minimum size=4pt},
  note/.style={font=\FigureNoteFont\fontsize{7.5}{9}\selectfont,text=FigureMuted,align=center}]
\begin{scope}
  \fill[FigureBlueFill] (0,0) rectangle (1,1);
  \draw[FigureGrid,line width=.5pt] (0,0) grid[step=.5] (1,1);
  \draw[FigureStroke,line width=.7pt] (0,0) rectangle (1,1);
  \draw[FigureGreen,line width=.8pt,dashed] (-.25,-.25) rectangle (.75,.75);
  \foreach \i in {1,2}{\foreach \j in {1,2}{\node[c] at ({(2*\i-1)/4},{(2*\j-1)/4}){};}}
  \node[note,text=FigureGreen!80!FigureInk,anchor=north] at (.25,-.27) {$B(\bm a,\tfrac12)$};
  \node at (.5,-.55) {(a) $\varepsilon=\tfrac12$，$4$个点};
\end{scope}
\begin{scope}[xshift=5.4cm]
  \fill[FigureBlueFill] (0,0) rectangle (1,1);
  \draw[FigureGrid,line width=.5pt] (0,0) grid[step=1/3] (1,1);
  \draw[FigureStroke,line width=.7pt] (0,0) rectangle (1,1);
  \draw[FigureGreen,line width=.8pt,dashed] ({1/2-1/4},{1/2-1/4}) rectangle ({1/2+1/4},{1/2+1/4});
  \foreach \i in {1,2,3}{\foreach \j in {1,2,3}{\node[c] at ({(2*\i-1)/6},{(2*\j-1)/6}){};}}
  \node at (.5,-.55) {(b) $\varepsilon=\tfrac14$，$9$个点};
\end{scope}
\end{tikzpicture}
